% !TeX root = Report.tex
\chapter{Численное сравнение по точности оценок параметров $\rho$ и $\theta$}

\section{Дизайн численного эксперимента}

Во всех опытах используется изображение размера $100\times100$, содержащее
одну прямую, и аддитивный гауссовский шум с уровнем $\sigma=0.2$. Для каждой
конфигурации выполняется $100$ повторов, причем для разных методов
предобработки используются одни и те же реализации шума. Шаг дискретизации
параметров преобразования Хафа фиксирован:
\[
h_\rho = 1,
\qquad
h_\theta = 0.01.
\]

Дополнительно тот же набор опытов был повторен для меньшего уровня шума
$\sigma=0.05$. Это позволяет сопоставить поведение методов не только в базовом
режиме, но и в более благоприятных условиях, когда влияние шумовой компоненты
ослаблено.

Рассматриваются четыре конфигурации прямой, приведенные в
табл.~\ref{tab:chapter3-configs}.

\begin{table}[H]
	\centering
	\caption{Конфигурации прямой в численном эксперименте}
	\label{tab:chapter3-configs}
	\resizebox{\textwidth}{!}{
	\begin{tabular}{|l|l|c|c|c|c|}
		\hline
		id & Конфигурация & $a$ & $b$ & $\rho$ & $\theta$ \\
		\hline
		\texttt{diag\_pos\_full} & Положительная диагональ & $1.00$ & $0.00$ & $0.00000$ & $2.35619$ \\
		\hline
		\texttt{diag\_neg\_full} & Отрицательная диагональ & $-1.00$ & $101.00$ & $71.41778$ & $0.78540$ \\
		\hline
		\texttt{steep\_full} & Крутая прямая & $2.00$ & $-1.00$ & $-0.44721$ & $2.67795$ \\
		\hline
		\texttt{steep\_shifted} & Крутая смещенная & $2.00$ & $-40.00$ & $-17.88854$ & $2.67795$ \\
		\hline
	\end{tabular}
	}
\end{table}

Каждая конфигурация из табл.~\ref{tab:chapter3-configs} задает не непрерывную
прямую как геометрическое множество, а ее дискретное представление на сетке
пикселей. Для фиксированных параметров $a$ и $b$ на изображении размера
\[
N_{\mathrm{row}} \times N_{\mathrm{col}}
\]
рассматривается прямая
\[
y = ax + b,
\]
а ее дискретное представление задается множеством пикселей
\[
\Gamma_{a,b}
=
\left\{
\bigl(x,[ax+b]\bigr)
\colon
x\in\{1,\dots,N_{\mathrm{col}}\},
\quad
1 \le [ax+b] \le N_{\mathrm{row}}
\right\}.
\]
Здесь $[\cdot]$ обозначает округление до ближайшего целого. Иными словами,
столбцовый индекс $x$ пробегает все целые значения от $1$ до $N_{\mathrm{col}}$,
а строковый индекс определяется как ближайшее целое к значению $ax+b$.
Тогда
\[
\Gamma_{a,b}
\subset
\{1,\dots,N_{\mathrm{col}}\}\times\{1,\dots,N_{\mathrm{row}}\},
\]
и матрица изображения задается правилом
\[
S^{(a,b)}_{y,x} =
\mathbf{1}\bigl\{(x,y)\in \Gamma_{a,b}\bigr\},
\]
а пиксели вне $\Gamma_{a,b}$ полагаются нулевыми.

Визуальный вид рассматриваемых конфигураций приведен на
рис.~\ref{fig:chapter3-config-preview}.

\begin{figure}[H]
	\centering
	\includegraphics[width=0.9\textwidth]{configuration-preview-1.pdf}
	\caption{Конфигурации прямой, используемые в численном эксперименте}
	\label{fig:chapter3-config-preview}
\end{figure}

\section{Сравниваемые методы и их параметры}

Активными, как и ранее, называются
пиксели с интенсивностью больше нуля после предобработки.

\begin{table}[H]
	\centering
	\caption{Методы предобработки и их параметры}
	\label{tab:chapter3-methods}
	\resizebox{\textwidth}{!}{
	\begin{tabular}{|l|l|c|}
		\hline
		Метод & Параметры & Среднее число активных пикселей \\
		\hline
		CSSA row.row & ранг аппроксимации num.line=1 & $3988.26$ \\
		\hline
		Median & окно $3\times3$ & $5128.47$ \\
		\hline
		Wiener & окно $5\times5$ & $9999.94$ \\
		\hline
		BM3D & параметр шума $\sigma_{psd}=0.2$ & $4229.38$ \\
		\hline
		Quantile 0.9 & порог по квантилю $q=0.9$ & $1000.00$ \\
		\hline
	\end{tabular}}
\end{table}

Как видно из таблицы ~\ref{tab:chapter3-methods} преобразование Хафа работает с множеством активных пикселей, поэтому для
дальнейшего сравнения существенна не только визуальная гладкость результата,
но и то, сколько точек остается на бинаризованном изображении. Любой шум, 
значения которого больше нуля, участвует в преобразовании. 
В этом смысле \texttt{Wiener} и \texttt{Median} оставляют существенно более плотное множество
активных пикселей, чем \texttt{Quantile 0.9}, а \texttt{CSSA row.row} и
\texttt{BM3D} занимают промежуточное положение.

\section{Точность оценивания при шуме $\sigma=0.2$}

Результаты для базового уровня шума $\sigma=0.2$ приведены в
табл.~\ref{tab:chapter3-rho-config},
\ref{tab:chapter3-theta-config},
\ref{tab:chapter3-rho-config-median}
и~\ref{tab:chapter3-theta-config-median}.
В каждом столбце жирным выделены наименьшие значения.

\begin{table}[H]
	\centering
	\caption{Средние ошибки оценивания параметра $\rho$ по конфигурациям прямой}
	\label{tab:chapter3-rho-config}
	\resizebox{\textwidth}{!}{
	\begin{tabular}{|l|c|c|c|c|}
		\hline
		Метод &
		\shortstack{Положительная\\диагональ} &
		\shortstack{Отрицательная\\диагональ} &
		\shortstack{Крутая\\прямая} &
		\shortstack{Крутая\\смещенная} \\
		\hline
		CSSA row.row & $\textbf{0.00000}$ & $0.20579$ & $1829.92372$ & $\textbf{0.01242}$ \\
		\hline
		Median & $\textbf{0.00000}$ & $0.21401$ & $101.94071$ & $65.96654$ \\
		\hline
		Wiener & $\textbf{0.00000}$ & $5100.50000$ & $\textbf{0.20000}$ & $320.00000$ \\
		\hline
		BM3D & $2.33000$ & $74.43233$ & $621.94095$ & $70.08796$ \\
		\hline
		Quantile 0.9 & $\textbf{0.00000}$ & $\textbf{0.17454}$ & $\textbf{0.20000}$ & $\textbf{0.01242}$ \\
		\hline
	\end{tabular}}
\end{table}

Табл.~\ref{tab:chapter3-rho-config} показывает, что основная проблема
\texttt{CSSA row.row} сосредоточена в конфигурации ``крутая прямая'', где
средняя ошибка по $\rho$ возрастает до $1829.92372$. Для диагональных и
смещенной крутой конфигураций метод остается существенно точнее. Таким
образом, ухудшение среднего качества для \texttt{CSSA row.row} вызвано не
равномерным снижением точности, а отдельной неустойчивой конфигурацией.

Для того же уровня шума сравнение по углу приведено в
табл.~\ref{tab:chapter3-theta-config}.

\begin{table}[H]
	\centering
	\caption{Средние ошибки оценивания параметра $\theta$ по конфигурациям прямой}
	\label{tab:chapter3-theta-config}
	\resizebox{\textwidth}{!}{
	\begin{tabular}{|l|c|c|c|c|}
		\hline
		Метод &
		\shortstack{Положительная\\диагональ} &
		\shortstack{Отрицательная\\диагональ} &
		\shortstack{Крутая\\прямая} &
		\shortstack{Крутая\\смещенная} \\
		\hline
		CSSA row.row & $\textbf{0.00001}$ & $0.00005$ & $2.20523$ & $\textbf{0.00000}$ \\
		\hline
		Median & $\textbf{0.00001}$ & $0.00003$ & $0.03602$ & $0.03045$ \\
		\hline
		Wiener & $\textbf{0.00001}$ & $2.45546$ & $0.10109$ & $0.10109$ \\
		\hline
		BM3D & $0.00006$ & $0.03748$ & $0.38626$ & $0.25210$ \\
		\hline
		Quantile 0.9 & $\textbf{0.00001}$ & $\textbf{0.00002}$ & $\textbf{0.00006}$ & $\textbf{0.00000}$ \\
		\hline
	\end{tabular}}
\end{table}

Наиболее устойчивым методом по параметру $\theta$ также оказывается
\texttt{Quantile 0.9}. Для \texttt{CSSA row.row} главный вклад в среднюю
ошибку снова дает конфигурация ``крутая прямая'', где ошибка по $\theta$
достигает $2.20523$. Следовательно, у данного метода сохраняется та же нестабильность, что и при сравнении по $\rho$.

Для интерпретации этой неустойчивости полезно отдельно рассмотреть
однострочный эксперимент. В строке длины $N=100$ фиксируется единственный
пиксель сигнала с индексом $j$, после чего к строке добавляется гауссовский шум. Ошибка
локализации определяется как $|\hat j-j|$, где $\hat j$ --- положение максимума
после восстановления. В стандартной для SSA нормированной записи частота такого
сигнала равна
\[
\omega_j=\frac{j-1}{N}.
\]
Далее рассматривается диапазон частот $0\le\omega_j\le1/2$. В эксперименте
использовались строки длины $N=100$, уровень шума
$\sigma=0.2$, а размер выборки для каждой точки графика составлял $1000$. Результат приведен на
рис.~\ref{fig:chapter3-frequency-localization}.

\begin{figure}[H]
	\centering
	\includegraphics[width=0.82\textwidth]{frequency_localization.pdf}
	\caption{Зависимость средней ошибки локализации максимума от нормированной частоты}
	\label{fig:chapter3-frequency-localization}
\end{figure}

Наблюдаемый частотный профиль следует понимать как эмпирическое свойство данной
однокомпонентной реконструкции при выбранном уровне шума, а не как общий
аналитический закон CSSA. Он показывает, что ошибки локализации максимума
сильнее проявляются в области малых частот.

Медианные ошибки для того же уровня шума приведены в
табл.~\ref{tab:chapter3-rho-config-median} и
табл.~\ref{tab:chapter3-theta-config-median}. В отличие от средних, медианы
характеризуют типичное поведение метода и существенно меньше чувствительны к
редким катастрофическим выбросам.

\begin{table}[H]
	\centering
	\caption{Медианные ошибки оценивания параметра $\rho$ по конфигурациям прямой}
	\label{tab:chapter3-rho-config-median}
	\resizebox{\textwidth}{!}{
	\begin{tabular}{|l|c|c|c|c|}
		\hline
		Метод &
		\shortstack{Положительная\\диагональ} &
		\shortstack{Отрицательная\\диагональ} &
		\shortstack{Крутая\\прямая} &
		\shortstack{Крутая\\смещенная} \\
		\hline
		CSSA row.row & $\textbf{0.00000}$ & $\textbf{0.17454}$ & $12.62229$ & $\textbf{0.01242}$ \\
		\hline
		Median & $\textbf{0.00000}$ & $\textbf{0.17454}$ & $0.30557$ & $\textbf{0.01242}$ \\
		\hline
		Wiener & $\textbf{0.00000}$ & $5100.50000$ & $\textbf{0.20000}$ & $320.00000$ \\
		\hline
		BM3D & $\textbf{0.00000}$ & $0.33897$ & $0.30557$ & $\textbf{0.01242}$ \\
		\hline
		Quantile 0.9 & $\textbf{0.00000}$ & $\textbf{0.17454}$ & $\textbf{0.20000}$ & $\textbf{0.01242}$ \\
		\hline
	\end{tabular}}
\end{table}

По медиане для параметра $\rho$ видно, что у \texttt{CSSA row.row} проблема
конфигурации ``крутая прямая'' действительно связана с выбросами: медианная
ошибка $12.62229$ намного меньше средней, хотя и остается заметно хуже, чем у
\texttt{Quantile 0.9} и \texttt{Wiener}. Следовательно, средняя ошибка здесь
раздувается редкими, но очень крупными промахами.

\begin{table}[H]
	\centering
	\caption{Медианные ошибки оценивания параметра $\theta$ по конфигурациям прямой}
	\label{tab:chapter3-theta-config-median}
	\resizebox{\textwidth}{!}{
	\begin{tabular}{|l|c|c|c|c|}
		\hline
		Метод &
		\shortstack{Положительная\\диагональ} &
		\shortstack{Отрицательная\\диагональ} &
		\shortstack{Крутая\\прямая} &
		\shortstack{Крутая\\смещенная} \\
		\hline
		CSSA row.row & $\textbf{0.000014}$ & $\textbf{0.000021}$ & $3.974631$ & $\textbf{0.000004}$ \\
		\hline
		Median & $\textbf{0.000014}$ & $\textbf{0.000021}$ & $\textbf{0.000004}$ & $\textbf{0.000004}$ \\
		\hline
		Wiener & $\textbf{0.000014}$ & $2.455460$ & $0.101089$ & $0.101089$ \\
		\hline
		BM3D & $\textbf{0.000014}$ & $0.000029$ & $0.000145$ & $\textbf{0.000004}$ \\
		\hline
		Quantile 0.9 & $\textbf{0.000014}$ & $\textbf{0.000021}$ & $0.000063$ & $\textbf{0.000004}$ \\
		\hline
	\end{tabular}}
\end{table}

Для параметра $\theta$ медианы дают еще более содержательную картину:
\texttt{CSSA row.row} на ``крутой прямой'' остается нестабильным не только по
среднему, но и по типичному поведению, тогда как \texttt{Median} и
\texttt{Quantile 0.9} дают малые медианные ошибки. Следовательно, в этой
конфигурации проблема \texttt{CSSA row.row} не сводится только к нескольким
редким выбросам.

\section{Точность оценивания при шуме $\sigma=0.05$}

Тот же численный эксперимент был повторен для уровня шума $\sigma=0.05$.
Средние ошибки оценивания параметров $\rho$ и $\theta$ по конфигурациям прямой
приведены в табл.~\ref{tab:chapter3-rho-config-low-noise} и
табл.~\ref{tab:chapter3-theta-config-low-noise}, а медианные ошибки ---
в табл.~\ref{tab:chapter3-rho-config-low-noise-median} и
\ref{tab:chapter3-theta-config-low-noise-median}. В каждом столбце жирным
выделены наименьшие значения.

\begin{table}[H]
	\centering
	\caption{Средние ошибки оценивания параметра $\rho$ по конфигурациям прямой при $\sigma=0.05$}
	\label{tab:chapter3-rho-config-low-noise}
	\resizebox{\textwidth}{!}{
	\begin{tabular}{|l|c|c|c|c|}
		\hline
		Метод &
		\shortstack{Положительная\\диагональ} &
		\shortstack{Отрицательная\\диагональ} &
		\shortstack{Крутая\\прямая} &
		\shortstack{Крутая\\смещенная} \\
		\hline
		CSSA row.row & $\textbf{0.00000}$ & $0.20743$ & $0.21795$ & $\textbf{0.01242}$ \\
		\hline
		Median & $\textbf{0.00000}$ & $0.22223$ & $0.53918$ & $1.94899$ \\
		\hline
		Wiener & $\textbf{0.00000}$ & $5100.50000$ & $\textbf{0.20000}$ & $319.32446$ \\
		\hline
		BM3D & $49.76000$ & $49.56193$ & $264.06563$ & $38.01611$ \\
		\hline
		Quantile 0.9 & $\textbf{0.00000}$ & $\textbf{0.17454}$ & $\textbf{0.20000}$ & $\textbf{0.01242}$ \\
		\hline
	\end{tabular}}
\end{table}

Для того же уровня шума сравнение по параметру $\theta$ приведено в
табл.~\ref{tab:chapter3-theta-config-low-noise}.

\begin{table}[H]
	\centering
	\caption{Средние ошибки оценивания параметра $\theta$ по конфигурациям прямой при $\sigma=0.05$}
	\label{tab:chapter3-theta-config-low-noise}
	\resizebox{\textwidth}{!}{
	\begin{tabular}{|l|c|c|c|c|}
		\hline
		Метод &
		\shortstack{Положительная\\диагональ} &
		\shortstack{Отрицательная\\диагональ} &
		\shortstack{Крутая\\прямая} &
		\shortstack{Крутая\\смещенная} \\
		\hline
		CSSA row.row & $\textbf{0.00001}$ & $0.00006$ & $0.00008$ & $\textbf{0.00000}$ \\
		\hline
		Median & $\textbf{0.00001}$ & $0.00003$ & $0.00205$ & $0.07917$ \\
		\hline
		Wiener & $\textbf{0.00001}$ & $2.45546$ & $0.10109$ & $0.10109$ \\
		\hline
		BM3D & $0.03227$ & $0.02471$ & $0.18344$ & $0.02539$ \\
		\hline
		Quantile 0.9 & $\textbf{0.00001}$ & $\textbf{0.00002}$ & $\textbf{0.00006}$ & $\textbf{0.00000}$ \\
		\hline
	\end{tabular}}
\end{table}

По сравнению с базовым случаем $\sigma=0.2$ при меньшем уровне шума резко
ослабевает нестабильность метода \texttt{CSSA row.row} на конфигурации
``крутая прямая'': средняя ошибка по $\rho$ уменьшается с $1829.92372$ до
$0.21795$, а средняя ошибка по $\theta$ --- с $2.20523$ до $0.00008$.
Следовательно, аномально большие ошибки этого метода связаны не столько с самой
конфигурацией прямой, сколько с сильным влиянием шума на результат. При этом \texttt{Quantile 0.9} сохраняет наилучшие или совместно
наилучшие значения во всех четырех конфигурациях, а \texttt{Wiener} даже при
меньшем шуме остается неустойчивым на отрицательной диагонали и смещенной
крутой прямой.

\begin{table}[H]
	\centering
	\caption{Медианные ошибки оценивания параметра $\rho$ по конфигурациям прямой при $\sigma=0.05$}
	\label{tab:chapter3-rho-config-low-noise-median}
	\resizebox{\textwidth}{!}{
	\begin{tabular}{|l|c|c|c|c|}
		\hline
		Метод &
		\shortstack{Положительная\\диагональ} &
		\shortstack{Отрицательная\\диагональ} &
		\shortstack{Крутая\\прямая} &
		\shortstack{Крутая\\смещенная} \\
		\hline
		CSSA row.row & $\textbf{0.00000}$ & $\textbf{0.17454}$ & $\textbf{0.20000}$ & $\textbf{0.01242}$ \\
		\hline
		Median & $\textbf{0.00000}$ & $\textbf{0.17454}$ & $0.30557$ & $\textbf{0.01242}$ \\
		\hline
		Wiener & $\textbf{0.00000}$ & $5100.50000$ & $\textbf{0.20000}$ & $320.00000$ \\
		\hline
		BM3D & $\textbf{0.00000}$ & $0.33897$ & $\textbf{0.20000}$ & $\textbf{0.01242}$ \\
		\hline
		Quantile 0.9 & $\textbf{0.00000}$ & $\textbf{0.17454}$ & $\textbf{0.20000}$ & $\textbf{0.01242}$ \\
		\hline
	\end{tabular}}
\end{table}

\begin{table}[H]
	\centering
	\caption{Медианные ошибки оценивания параметра $\theta$ по конфигурациям прямой при $\sigma=0.05$}
	\label{tab:chapter3-theta-config-low-noise-median}
	\resizebox{\textwidth}{!}{
	\begin{tabular}{|l|c|c|c|c|}
		\hline
		Метод &
		\shortstack{Положительная\\диагональ} &
		\shortstack{Отрицательная\\диагональ} &
		\shortstack{Крутая\\прямая} &
		\shortstack{Крутая\\смещенная} \\
		\hline
		CSSA row.row & $\textbf{0.000014}$ & $\textbf{0.000021}$ & $0.000063$ & $\textbf{0.000004}$ \\
		\hline
		Median & $\textbf{0.000014}$ & $\textbf{0.000021}$ & $\textbf{0.000004}$ & $\textbf{0.000004}$ \\
		\hline
		Wiener & $\textbf{0.000014}$ & $2.455460$ & $0.101089$ & $0.101089$ \\
		\hline
		BM3D & $\textbf{0.000014}$ & $0.000029$ & $0.000063$ & $\textbf{0.000004}$ \\
		\hline
		Quantile 0.9 & $\textbf{0.000014}$ & $\textbf{0.000021}$ & $0.000063$ & $\textbf{0.000004}$ \\
		\hline
	\end{tabular}}
\end{table}

Медианы при $\sigma=0.05$ подтверждают вывод о шумовой природе нестабильности
\texttt{CSSA row.row}. На конфигурации ``крутая прямая'' медианная ошибка по
$\rho$ уже совпадает с наилучшим значением $0.20000$, а медианная ошибка по
$\theta$ уменьшается до $0.000063$. Следовательно, при меньшем уровне шума
типичное поведение \texttt{CSSA row.row} становится сопоставимым с
\texttt{Quantile 0.9}, тогда как заметные проблемы сохраняются в основном у
\texttt{Wiener}.

\section{Влияние шага дискретизации параметров Хафа}

Дополнительный опыт был проведен для одной прямой
\[
y = 2x - 1
\]
на изображении размера $100\times100$ при уровне шума $\sigma=0.05$. Во всех
запусках использовалась предобработка \texttt{CSSA row.row}, а затем
преобразование Хафа применялось с разными шагами дискретизации по $\rho$ и
$\theta$. Для каждой сетки выполнялось $100$ повторов.

\begin{table}[H]
	\centering
	\caption{Влияние шага дискретизации параметров Хафа на точность оценивания после \texttt{CSSA row.row}}
	\label{tab:chapter3-discretization-step}
	\resizebox{\textwidth}{!}{
	\begin{tabular}{|l|c|c|c|c|c|c|}
\hline
Сетка & $h_\rho$ & $h_\theta$ & Средняя $\Delta_\rho$ & Медианная $\Delta_\rho$ & Средняя $\Delta_\theta$ & Медианная $\Delta_\theta$ \\
\hline
Крупный шаг & 5.0000 & 1.0000 & 899.246070 & 995.578330 & 0.103719 & 0.103719 \\
\hline
Средний шаг & 1.0000 & 0.0100 & 0.200000 & 0.200000 & 0.000063 & 0.000063 \\
\hline
Мелкий шаг & 0.5000 & 0.0050 & 0.200000 & 0.200000 & 0.000009 & 0.000009 \\
\hline
Очень мелкий шаг & 0.2500 & 0.0025 & \textbf{0.002786} & \textbf{0.002786} & \textbf{0.000000} & \textbf{0.000000} \\
\hline
\end{tabular}

	}
\end{table}

Табл.~\ref{tab:chapter3-discretization-step} показывает, что переход от
крупного шага $(h_\rho, h_\theta)=(5, 1)$ к сетке $(1, 0.01)$ резко уменьшает
ошибки по обоим параметрам. Дальнейшее уменьшение шага по $\theta$ продолжает
снижать ошибку по углу, тогда как ошибка по $\rho$ заметно уменьшается только
на самой мелкой сетке. Это связано с тем, что при более грубой дискретизации
максимум аккумуляторной матрицы может быть локализован лишь в относительно
крупной окрестности истинного значения. Следовательно, уменьшение шага
дискретизации параметров преобразования Хафа приводит к более точной оценке
параметров прямой, хотя выигрыш по $\rho$ и по $\theta$ проявляется
неодновременно.

\section{Итог сравнения}

Численные результаты хорошо согласуются с наблюдениями главы~2. Во-первых,
для метода Хафа критична не визуальная гладкость восстановленного
изображения, а то, как расположены множества активных пикселей.
Именно поэтому \texttt{Wiener}, который в среднем оставляет практически все пискели 
активными, дает наихудшие оценки: преобразование Хафа в этом случае
получает почти сплошное множество точек и теряет устойчивость.

Во-вторых, результаты \texttt{CSSA row.row} подтверждают выводы
рис.~\ref{fig:row-row-components} и рис.~\ref{fig:one-row-bad-case}. Этот
метод дает малые ошибки в диагональных и смещенной крутой конфигурациях,
однако для конфигурации ``крутая прямая'' возникают эффекты растекания и
смещения максимума вследствие неправильного выбора компоненты. Для
преобразования Хафа такие сбои существенны, поскольку даже небольшое число
ложных активных пикселей может изменить положение максимума в пространстве
параметров.

Дополнительный прогон при $\sigma=0.05$ показывает, что эта неустойчивость
носит ярко выраженный шумовой характер. При меньшем уровне шума
\texttt{CSSA row.row} практически сравнивается с \texttt{Quantile 0.9} по
точности оценивания параметров прямой и перестает давать катастрофические
промахи на конфигурации ``крутая прямая''.

В-третьих, устойчивость \texttt{Quantile 0.9} связана с тем, что этот метод
не только подавляет слабые значения, но и жестко ограничивает множество
активных пикселей. В рассматриваемом эксперименте после такой предобработки
активными остаются в среднем $1000$ пикселей, тогда как после
\texttt{CSSA row.row} их около $4000$, после \texttt{Median} более $5000$, а
после \texttt{Wiener} почти $10000$. Для последующего применения
преобразования Хафа такая разница принципиальна.

Наконец, результаты \texttt{BM3D} и \texttt{Median} показывают, что хорошее
визуальное шумоподавление еще не гарантирует наилучшую оценку параметров
прямой. Эти методы подавляют шум заметно лучше, чем отсутствие обработки, но
после бинаризации все еще оставляют достаточно плотное множество активных
пикселей. Поэтому по точности оценок $\rho$ и $\theta$ они уступают
\texttt{Quantile 0.9}, а \texttt{CSSA row.row} превосходят лишь в тех
конфигурациях, где у построчной однокомпонентной реконструкции возникают
крупные сбои.
